% GENERATED FILE. Edit canonical lessons, then run npm run generate:books.
% canonical-source-hash: fnv1a64:4e03dd9c49d34d3e
% canonical-lessons: AR-C535-rafiq, AR-C535-aduww, AR-C535-khiyana, AR-C535-takhasama, AR-C535-hajara, AR-R535-review-528-535

\chapter{Companions and quarrels}
\label{ch:ar-companions-and-quarrels}
\hlchaptermodality{535}

\begin{chapteropening}
\textbf{By the end of this chapter:} \emph{I can talk about a companion, an enemy and betrayal, and say that people fell out or that someone left a place.}

Complete the last lesson of chapter 535: I can talk about a companion, an enemy and betrayal, and say that people fell out or that someone left a place.
\end{chapteropening}

\section[\texorpdfstring{\ar{رفيق} (rafīq)}{رفيق (rafīq)}]{\ar{رفيق} (rafīq) — a companion}

\label{lesson:AR-C535-rafiq}



\begin{quote}
Before the new one: say the Arabic for esteem, then the Arabic for to pardon.
\end{quote}

\subsection*{You'll want to know: \ar{رفيق}}
\textbf{\ar{رفيق}} — \emph{rafīq} — \textquotedblleft{}a companion\textquotedblright{}.

\begin{quote}
\ar{كان} \ar{الكلب} \ar{رفيق} \ar{جدي} \ar{في} \ar{الجبل}.
\end{quote}

\begin{quote}
\emph{kāna al-kalb rafīq jaddī fī al-jabal.}
\end{quote}

\begin{quote}
The dog was my grandfather's companion in the mountains.
\end{quote}

Plural: \ar{رفاق}, \emph{rifāq}.

\begin{cousinweb}
\textbf{\ar{رفيق}} (\emph{rafīq}) is built on the native Arabic root \emph{r-f-q}, \textquotedblleft{}to be gentle; to be a companion\textquotedblright{} (Lane). A rafīq goes along with you.
\end{cousinweb}

\begin{grammarlens}[title={What you've built}]
One more word for describing what you lived and felt.
\end{grammarlens}

\subsection*{Your turn}
\begin{itemize}
\raggedright
  \item \emph{Say it:} \emph{rafīq}
  \item \emph{Say it:} \emph{rafīq}, once more
  \item \emph{Say it:} \emph{kāna al-kalb rafīq jaddī fī al-jabal.} — The dog was my grandfather's companion in the mountains.
  \item \emph{Recall:} say the Arabic for esteem, then the Arabic for to pardon, then say \emph{rafīq} again
\end{itemize}

\begin{tcolorbox}[breakable,colback=teal!4,colframe=teal!35!black,title={Before you move on}]
What does \ar{رفيق} mean? (\textquotedblleft{}A companion\textquotedblright{}.) Say it once more.
\end{tcolorbox}
\section[\texorpdfstring{\ar{عدو} (ʿaduww)}{عدو (ʿaduww)}]{\ar{عدو} (ʿaduww) — an enemy}

\label{lesson:AR-C535-aduww}



\begin{quote}
Before the new one: say the Arabic for to pardon, then the Arabic for a companion.
\end{quote}

\subsection*{You'll want to know: \ar{عدو}}
\textbf{\ar{عدو}} — \emph{ʿaduww} — \textquotedblleft{}an enemy\textquotedblright{}.

\begin{quote}
\ar{القط} \ar{والفأر} \ar{عدوان} \ar{في} \ar{القصص}.
\end{quote}

\begin{quote}
\emph{al-qiṭṭ wa-l-faʾr ʿaduwwān fī al-qiṣaṣ.}
\end{quote}

\begin{quote}
The cat and the mouse are enemies in stories.
\end{quote}

Plural: \ar{أعداء}, \emph{aʿdāʾ}.

\begin{cousinweb}
\textbf{\ar{عدو}} (\emph{ʿaduww}) is built on the native Arabic root \emph{ʿ-d-w}, \textquotedblleft{}to run, to pass beyond\textquotedblright{} (Lane). An ʿaduww is one who runs against you, who goes beyond what is right.
\end{cousinweb}

\begin{grammarlens}[title={What you've built}]
One more word for describing what you lived and felt.
\end{grammarlens}

\subsection*{Your turn}
\begin{itemize}
\raggedright
  \item \emph{Say it:} \emph{ʿaduww}
  \item \emph{Say it:} \emph{ʿaduww}, once more
  \item \emph{Say it:} \emph{al-qiṭṭ wa-l-faʾr ʿaduwwān fī al-qiṣaṣ.} — The cat and the mouse are enemies in stories.
  \item \emph{Recall:} say the Arabic for to pardon, then the Arabic for a companion, then say \emph{ʿaduww} again
\end{itemize}

\begin{tcolorbox}[breakable,colback=teal!4,colframe=teal!35!black,title={Before you move on}]
What does \ar{عدو} mean? (\textquotedblleft{}An enemy\textquotedblright{}.) Say it once more.
\end{tcolorbox}
\section[\texorpdfstring{\ar{خيانة} (khiyāna)}{خيانة (khiyāna)}]{\ar{خيانة} (khiyāna) — betrayal}

\label{lesson:AR-C535-khiyana}



\begin{quote}
Before the new one: say the Arabic for a companion, then the Arabic for an enemy.
\end{quote}

\subsection*{You'll want to know: \ar{خيانة}}
\textbf{\ar{خيانة}} — \emph{khiyāna} — \textquotedblleft{}betrayal\textquotedblright{}.

\begin{quote}
\ar{الكذب} \ar{على} \ar{صديق} \ar{خيانة}.
\end{quote}

\begin{quote}
\emph{al-kadhib ʿalā ṣadīq khiyāna.}
\end{quote}

\begin{quote}
Lying to a friend is a betrayal.
\end{quote}

\begin{cousinweb}
\textbf{\ar{خيانة}} (\emph{khiyāna}) is built on the native Arabic root \emph{kh-w-n}, \textquotedblleft{}to be false, to betray\textquotedblright{} (Lane). Khiyāna is the verbal noun of khāna.
\end{cousinweb}

\begin{grammarlens}[title={What you've built}]
One more word for describing what you lived and felt.
\end{grammarlens}

\subsection*{Your turn}
\begin{itemize}
\raggedright
  \item \emph{Say it:} \emph{khiyāna}
  \item \emph{Say it:} \emph{khiyāna}, once more
  \item \emph{Say it:} \emph{al-kadhib ʿalā ṣadīq khiyāna.} — Lying to a friend is a betrayal.
  \item \emph{Recall:} say the Arabic for a companion, then the Arabic for an enemy, then say \emph{khiyāna} again
\end{itemize}

\begin{tcolorbox}[breakable,colback=teal!4,colframe=teal!35!black,title={Before you move on}]
What does \ar{خيانة} mean? (\textquotedblleft{}Betrayal\textquotedblright{}.) Say it once more.
\end{tcolorbox}
\section[\texorpdfstring{\ar{تخاصم} (takhāṣama)}{تخاصم (takhāṣama)}]{\ar{تخاصم} (takhāṣama) — to fall out (with each other)}

\label{lesson:AR-C535-takhasama}



\begin{quote}
Before the new one: say the Arabic for an enemy, then the Arabic for betrayal.
\end{quote}

\subsection*{You'll want to know: \ar{تخاصم}}
\textbf{\ar{تخاصم}} — \emph{takhāṣama} — \textquotedblleft{}to fall out (with each other)\textquotedblright{}.

\begin{quote}
\ar{تخاصم} \ar{الأخوان} \ar{على} \ar{المال}.
\end{quote}

\begin{quote}
\emph{takhāṣama al-akhawān ʿalā al-māl.}
\end{quote}

\begin{quote}
The two brothers fell out over money.
\end{quote}

Present: \ar{يتخاصم}, \emph{yatakhāṣam}.

\begin{cousinweb}
\textbf{\ar{تخاصم}} (\emph{takhāṣama}) is built on the native Arabic root \emph{kh-ṣ-m}, \textquotedblleft{}to dispute, to contend\textquotedblright{} (Lane). Form VI takhāṣama is to dispute with one another.
\end{cousinweb}

\begin{grammarlens}[title={What you've built}]
One more word for describing what you lived and felt.
\end{grammarlens}

\subsection*{Your turn}
\begin{itemize}
\raggedright
  \item \emph{Say it:} \emph{takhāṣama}
  \item \emph{Say it:} \emph{takhāṣama}, once more
  \item \emph{Say it:} \emph{takhāṣama al-akhawān ʿalā al-māl.} — The two brothers fell out over money.
  \item \emph{Recall:} say the Arabic for an enemy, then the Arabic for betrayal, then say \emph{takhāṣama} again
\end{itemize}

\begin{tcolorbox}[breakable,colback=teal!4,colframe=teal!35!black,title={Before you move on}]
What does \ar{تخاصم} mean? (\textquotedblleft{}To fall out (with each other)\textquotedblright{}.) Say it once more.
\end{tcolorbox}
\section[\texorpdfstring{\ar{هجر} (hajara)}{هجر (hajara)}]{\ar{هجر} (hajara) — to abandon, to leave}

\label{lesson:AR-C535-hajara}



\begin{quote}
Before the new one: say the Arabic for betrayal, then the Arabic for to fall out.
\end{quote}

\subsection*{You'll want to know: \ar{هجر}}
\textbf{\ar{هجر}} — \emph{hajara} — \textquotedblleft{}to abandon, to leave\textquotedblright{}.

\begin{quote}
\ar{هجر} \ar{الشاب} \ar{قريته} \ar{وسكن} \ar{في} \ar{المدينة}.
\end{quote}

\begin{quote}
\emph{hajara ash-shābb qaryatahu wa-sakana fī al-madīna.}
\end{quote}

\begin{quote}
The young man left his village and lived in the city.
\end{quote}

Present: \ar{يهجر}, \emph{yahjur}.

\begin{cousinweb}
\textbf{\ar{هجر}} (\emph{hajara}) is built on the native Arabic root \emph{h-j-r}, \textquotedblleft{}to cut off, to forsake\textquotedblright{} (Lane). Lane gives hajara, aor. yahjuru, as to cut oneself off from someone.
\end{cousinweb}

\begin{grammarlens}[title={What you've built}]
One more word for describing what you lived and felt.
\end{grammarlens}

\subsection*{Your turn}
\begin{itemize}
\raggedright
  \item \emph{Say it:} \emph{hajara}
  \item \emph{Say it:} \emph{hajara}, once more
  \item \emph{Say it:} \emph{hajara ash-shābb qaryatahu wa-sakana fī al-madīna.} — The young man left his village and lived in the city.
  \item \emph{Recall:} say the Arabic for betrayal, then the Arabic for to fall out, then say \emph{hajara} again
\end{itemize}

\begin{tcolorbox}[breakable,colback=teal!4,colframe=teal!35!black,title={Before you move on}]
What does \ar{هجر} mean? (\textquotedblleft{}To abandon, to leave\textquotedblright{}.) Say it once more.
\end{tcolorbox}
\section[(dialogue)]{Words from chapters 528–535 — a review}

\label{lesson:AR-R535-review-528-535}



\begin{quote}
Say the Arabic for a fiancé and marriage.
\end{quote}

\subsection*{Your turn}
Say these aloud:

\begin{itemize}
\raggedright
  \item a fiancé, marriage, a mother-in-law, single, to love passionately
  \item a son-in-law, a daughter-in-law, an in-law, a brother-in-law, a marriage break-up
  \item pregnant, a cradle, to breastfeed, to wean, to bring up
  \item a teenager, to get older, to ripen, old age, a nanny
  \item a parent, a stepson, kinship, to lead a life, to die
  \item to look after, to pamper, to cradle, to stand by, a loved one
  \item to obey, to disobey, to scold, esteem, to pardon
  \item a companion, an enemy, betrayal, to fall out, to abandon
\end{itemize}

\begin{tcolorbox}[breakable,colback=teal!4,colframe=teal!35!black,title={Before you move on}]
A fiancé? (\textbf{\ar{خطيب}}, \emph{khaṭīb}.) Marriage? (\textbf{\ar{زواج}}, \emph{zawāj}.) A mother-in-law? (\textbf{\ar{حماة}}, \emph{ḥamāh}.) Single? (\textbf{\ar{أعزب}}, \emph{aʿzab}.) To love passionately? (\textbf{\ar{عشق}}, \emph{ʿashiqa}.) A son-in-law? (\textbf{\ar{صهر}}, \emph{ṣihr}.) A daughter-in-law? (\textbf{\ar{كنة}}, \emph{kanna}.) An in-law? (\textbf{\ar{نسيب}}, \emph{nasīb}.) A brother-in-law? (\textbf{\ar{عديل}}, \emph{ʿadīl}.) A marriage break-up? (\textbf{\ar{طلاق}}, \emph{ṭalāq}.) Pregnant? (\textbf{\ar{حامل}}, \emph{ḥāmil}.) A cradle? (\textbf{\ar{مهد}}, \emph{mahd}.) To breastfeed? (\textbf{\ar{أرضع}}, \emph{arḍaʿa}.) To wean? (\textbf{\ar{فطم}}, \emph{faṭama}.) To bring up? (\textbf{\ar{ربى}}, \emph{rabbā}.) A teenager? (\textbf{\ar{مراهق}}, \emph{murāhiq}.) To get older? (\textbf{\ar{كبر}}, \emph{kabira}.) To ripen? (\textbf{\ar{نضج}}, \emph{naḍija}.) Old age? (\textbf{\ar{شيخوخة}}, \emph{shaykhūkha}.) A nanny? (\textbf{\ar{مربية}}, \emph{murabbiya}.) A parent? (\textbf{\ar{والد}}, \emph{wālid}.) A stepson? (\textbf{\ar{ربيب}}, \emph{rabīb}.) Kinship? (\textbf{\ar{قرابة}}, \emph{qarāba}.) To lead a life? (\textbf{\ar{عاش}}, \emph{ʿāsha}.) To die? (\textbf{\ar{مات}}, \emph{māta}.) To look after? (\textbf{\ar{اعتنى}}, \emph{iʿtanā}.) To pamper? (\textbf{\ar{دلل}}, \emph{dallala}.) To cradle? (\textbf{\ar{حضن}}, \emph{ḥaḍana}.) To stand by? (\textbf{\ar{ساند}}, \emph{sānada}.) A loved one? (\textbf{\ar{حبيب}}, \emph{ḥabīb}.) To obey? (\textbf{\ar{أطاع}}, \emph{aṭāʿa}.) To disobey? (\textbf{\ar{عصى}}, \emph{ʿaṣā}.) To scold? (\textbf{\ar{وبخ}}, \emph{wabbakha}.) Esteem? (\textbf{\ar{احترام}}, \emph{iḥtirām}.) To pardon? (\textbf{\ar{عفا}}, \emph{ʿafā}.) A companion? (\textbf{\ar{رفيق}}, \emph{rafīq}.) An enemy? (\textbf{\ar{عدو}}, \emph{ʿaduww}.) Betrayal? (\textbf{\ar{خيانة}}, \emph{khiyāna}.) To fall out? (\textbf{\ar{تخاصم}}, \emph{takhāṣama}.) To abandon? (\textbf{\ar{هجر}}, \emph{hajara}.)
\end{tcolorbox}
